% Source: log-softmax_trick.tex
\begin{notechapter}
\chapter{The Softmax Function and its Numerical Stability}\label{ch:log-softmax}
\chaptersubtitle{The log-sum-exp trick}

\section{Introduction}
The softmax function is a fundamental component of many machine learning 
models, particularly in classification tasks. It converts a vector of raw 
scores (logits) into a probability distribution. However, computing the 
softmax directly can lead to numerical instability. This document provides 
a detailed explanation of the softmax function, its challenges, and how to 
stabilize it using the log-softmax trick. Additionally, it explores the 
connection between softmax and cross-entropy loss, along with practical 
implementations in deep learning frameworks.

\section{The Softmax Function}
The softmax function takes a vector of logits \( x = [x_1, x_2, \dots, x_n] \) 
and converts it into a probability distribution \( p = [p_1, p_2, \dots, p_n] \), 
where:
\[
p_i = \text{softmax}(x_i) = \frac{e^{x_i}}{\sum_{j=1}^n e^{x_j}}
\]
The output \( p_i \) represents the probability of the \( i \)-th class. The 
softmax function ensures that:
\begin{itemize}
    \item All outputs are non-negative: \( p_i \geq 0 \).
    \item The outputs sum to 1: \( \sum_{i=1}^n p_i = 1 \).
\end{itemize}

\section{Numerical Instability in Softmax}
Direct computation of the softmax can lead to numerical instability due to:
\begin{itemize}
    \item \textbf{Overflow:} If \( x_i \) is very large, \( e^{x_i} \) can 
    exceed the range of floating-point numbers, resulting in infinity.
    \item \textbf{Underflow:} If \( x_i \) is very small, \( e^{x_i} \) can 
    round to zero, leading to division by zero or undefined behavior.
\end{itemize}
For example, consider the logits \( x = [1000, 1001, 1002] \). Directly 
computing \( e^{1000} \) would result in overflow.

\section{Stabilizing Softmax with the Log-Sum-Exp Trick}
To address these issues, we use the \textbf{log-sum-exp trick}, which 
stabilizes the computation by subtracting the maximum logit \( c = \max(x_1, x_2, \dots, x_n) \).

\subsection{Stabilized Softmax}
The stabilized softmax is computed as:
\[
\text{softmax}(x_i) = \frac{e^{x_i - c}}{\sum_{j=1}^n e^{x_j - c}}
\]
Subtracting \( c \) ensures that the largest exponent is \( e^{c - c} = 1 \), 
and all other exponents are \( e^{x_i - c} \leq 1 \), preventing overflow.
Very negative \( x_i - c \) can still underflow to \( 0 \), but this is harmless:
the denominator is at least \( 1 \), so there is no division by zero, and the
corresponding probabilities are genuinely negligible.

\subsection{Log-Softmax}
For additional numerical stability, we often compute the \textbf{log-softmax} 
directly:
\[
\log(\text{softmax}(x_i)) = x_i - c - \log\left(\sum_{j=1}^n e^{x_j - c}\right)
\]
This avoids taking the logarithm of a probability that has underflowed to
\( 0 \) (which would give \( -\infty \)): the log-probability \( x_i - c - \log\sum_j e^{x_j - c} \)
remains finite and accurate even when \( p_i \) itself is not representable. It is
also cheaper and better conditioned for gradient computations.

\subsection{Why the Log-Sum-Exp Trick Works}
The log-sum-exp trick works because:
\begin{itemize}
    \item Subtracting \( c \) ensures that the largest exponent is 1, 
    preventing overflow.
        \item All other exponents are bounded between 0 and 1; some may
    underflow to 0, but the sum is at least 1, so the division and the
    logarithm of the sum are always well defined.
    \item The final result is mathematically equivalent to the original 
    softmax, as the subtraction cancels out in the numerator and denominator.
\end{itemize}

\section{Example: Stabilizing Softmax}
Consider the logits \( x = [1000, 1001, 1002] \):
\begin{enumerate}
    \item Compute \( c = \max(1000, 1001, 1002) = 1002 \).
    \item Subtract \( c \) from each logit:
    \[
    x' = [1000 - 1002, 1001 - 1002, 1002 - 1002] = [-2, -1, 0]
    \]
    \item Compute the stabilized softmax:
    \[
    \text{softmax}(x_i) = \frac{e^{x_i'}}{e^{-2} + e^{-1} + e^{0}}
    \]
    \[
    = \frac{e^{x_i'}}{0.1353 + 0.3679 + 1} = \frac{e^{x_i'}}{1.5032}
    \]
    \item Compute the probabilities:
    \[
    \text{softmax}(x_1) = \frac{e^{-2}}{1.5032} \approx 0.0900
    \]
    \[
    \text{softmax}(x_2) = \frac{e^{-1}}{1.5032} \approx 0.2447
    \]
    \[
    \text{softmax}(x_3) = \frac{e^{0}}{1.5032} \approx 0.6652
    \]
\end{enumerate}

\section{Cross-Entropy Loss with Log-Softmax}
The cross-entropy loss measures the difference between the true label 
distribution \( y_i \) and the predicted distribution \( p_i \):
\[
\text{Cross-Entropy} = -\sum_{i=1}^n y_i \log(p_i)
\]
When using log-softmax, the loss is computed as:
\[
\text{Cross-Entropy} = -\sum_{i=1}^n y_i \cdot \log(\text{softmax}(x_i))
\]
Substituting the log-softmax:
\[
\text{Cross-Entropy} = -x_k + c + \log\left(\sum_{j=1}^n e^{x_j - c}\right)
\]
where \( k \) is the index of the true class.

\section{Practical Implementation}
Deep learning frameworks like PyTorch and TensorFlow provide built-in 
functions for log-softmax and cross-entropy loss.

\subsection{PyTorch Example}
\begin{verbatim}
import torch
import torch.nn as nn

logits = torch.tensor([[1000.0, 1001.0, 1002.0]])
true_labels = torch.tensor([2])  # True class index

criterion = nn.CrossEntropyLoss()
loss = criterion(logits, true_labels)
print(loss.item())  # Output: 0.4076
\end{verbatim}

\subsection{TensorFlow Example}
\begin{verbatim}
import tensorflow as tf

logits = tf.constant([[1000.0, 1001.0, 1002.0]])
true_labels = tf.constant([2])  # True class index

loss = tf.nn.sparse_softmax_cross_entropy_with_logits(
    labels=true_labels, logits=logits
)
print(loss.numpy())  # Output: [0.4076]
\end{verbatim}

\section{Conclusion}
The softmax function is a key component of classification models, but its 
direct computation can lead to numerical instability. The log-softmax trick 
stabilizes the computation and is widely used in practice. Combined with 
cross-entropy loss, it forms the foundation of modern classification tasks 
in machine learning. By understanding and applying these concepts, you can 
build more robust and efficient machine learning models.
\end{notechapter}
